\documentclass[11pt,a4paper]{article}
\usepackage[hyperref]{acl}
\usepackage{times}
\usepackage{latexsym}
\usepackage{graphicx}
\usepackage{amsmath}
\usepackage{amssymb}
\usepackage{booktabs}
\usepackage{multirow}
\usepackage{xcolor}
\usepackage{tablefootnote}

\title{The Formulation Design Space for Neural Code Verification:\\A Systematic Comparison of Classification, Generation, and Trajectory Approaches}

\author{Anonymous}

\begin{document}
\maketitle

%==============================================================================
\begin{abstract}
Neural code verification---predicting whether a candidate patch passes a test suite without execution---can be formulated as classification (predict binary pass/fail), constrained generation (predict execution output then match), or trajectory scoring (predict overall resolution status). Which formulation should practitioners choose? We present the first controlled study spanning 3 formulations $\times$ 2 model families (Qwen, DeepSeek-Coder) $\times$ 2 complexity levels (function-level, repository-level), trained on a new per-test dataset derived from SWE-bench Verified (489 instances, 9,784 test-level samples). Our findings challenge the assumption that any single formulation dominates: formulation ranking is \emph{model-contingent}---classification outperforms generation on Qwen at repository scale (+24pp) due to tokenizer-specific output collapse, yet generation \emph{matches or exceeds} classification on DeepSeek-Coder where unparseable rates drop from 93\% to 12\%. Across all settings, classification and generation exhibit complementary error patterns (Cohen's $\kappa$ = 0.24--0.31), enabling oracle ensembles that lift accuracy by +2--3pp over the best single formulation. When accuracy is equivalent, classification offers 6--9$\times$ inference speedup via single-token decoding. We release our dataset and provide a practical decision framework for formulation selection conditioned on model family and task complexity.
\end{abstract}

%==============================================================================
\section{Introduction}

Large language models have dramatically improved automated code repair, with systems like SWE-agent achieving non-trivial resolve rates on real-world software engineering tasks \cite{jimenez2024swebench}. A critical bottleneck in test-time scaling of these systems is \emph{verification}: determining whether a candidate patch actually fixes the intended bug. Full test execution provides ground truth but costs \$1--3 and minutes per candidate, making it prohibitive when ranking hundreds of candidates.

Execution-free verification has emerged as a practical alternative. Existing approaches span a design space from trajectory-level binary scoring (SWE-RM; ``is this patch resolved?'') to fine-grained execution simulation (Code World Models; ``what output does this test produce?''). Yet no prior work systematically compares these formulations under controlled conditions---same base model, same data, same compute budget---across model families and complexity levels.

We address this gap with three contributions:

\begin{enumerate}
\item \textbf{Error Complementarity as Verification Strategy.} We demonstrate that classification and generation formulations exhibit systematically different failure modes---classification tends toward missed bugs (false negatives) while generation tends toward false alarms (false positives)---with only 15\% error type overlap. Oracle ensembles exploiting this complementarity achieve 89--91\% accuracy across both model families, lifting +2--3pp over the best single formulation.

\item \textbf{First Systematic Formulation Comparison.} We provide the first 3-formulation $\times$ 2-model-family $\times$ 2-complexity controlled comparison for neural code verification, revealing that formulation effectiveness is \emph{model-contingent}: no universal winner exists, and the interaction between model architecture and formulation choice is statistically significant ($p=0.03$).

\item \textbf{Efficiency--Accuracy Tradeoff Characterization.} When classification and generation achieve equivalent accuracy (as on DeepSeek-Coder: 86.85\% vs 86.83\%), classification provides 6--9$\times$ inference speedup through single-token decoding, establishing it as the Pareto-optimal choice in accuracy-equivalent regimes.
\end{enumerate}

Our findings have immediate practical implications: practitioners should not assume a single ``best'' formulation but rather select based on their model family and deploy formulation-diverse ensembles when verification reliability is critical.

%==============================================================================
\section{Related Work}

\paragraph{Neural Code Verification.}
Execution prediction for code has been explored through Code World Models \cite{gu2024cruxeval}, which predict stdout given code and input. SWE-RM \cite{zhang2025swerm} trains trajectory-level reward models for patch ranking. Our work differs by systematically comparing formulation choices rather than proposing a single new method.

\paragraph{Classification vs Generation in NLP.}
The tradeoff between discriminative classification and generative approaches is well-studied in NLP \cite{ng2001discriminative}. In code verification, this manifests as a choice between predicting a binary label (pass/fail) versus generating execution output and matching. We provide controlled empirical evidence for this tradeoff in the code domain.

\paragraph{Test-Time Scaling for Code Repair.}
Recent work on scaling test-time compute for code generation \cite{chen2024codegen} relies on verification to rank candidates. Our per-test formulation provides diagnostic granularity---identifying \emph{which} tests fail---potentially enabling targeted re-repair, though our downstream experiment shows this utility remains unrealized with current repair agents.

%==============================================================================
\section{Method}

\subsection{Task Formulation}

Given a code context $C$ (function or repository with patch applied) and a test case $t$, we study three verification formulations:

\paragraph{Classification (CLS).} Predict a single token $y \in \{\texttt{pass}, \texttt{fail}\}$:
$$P_\text{CLS}(y \mid C, t) = \text{softmax}(W h_\text{[EOS]})$$
where $h_\text{[EOS]}$ is the final hidden state after processing the input.

\paragraph{Constrained Word Model (CWM).} Generate the expected execution output $\hat{o}$, then determine pass/fail by string matching against the expected output:
$$\hat{o} = \arg\max \prod_{i=1}^{N} P(o_i \mid o_{<i}, C, t)$$
$$y_\text{CWM} = \mathbb{1}[\text{match}(\hat{o}, o^*)]$$

\paragraph{Trajectory (TRAJ).} Predict a single instance-level label $y \in \{\texttt{resolved}, \texttt{unresolved}\}$ aggregating all tests:
$$P_\text{TRAJ}(y \mid C, T) \text{ where } T = \{t_1, \ldots, t_k\}$$

\subsection{Dataset: SWE-bench-PerTest}

We construct a per-test evaluation dataset from SWE-bench Verified by re-executing gold patches in Docker containers to extract individual test outcomes (Table~\ref{tab:dataset}).

\begin{table}[t]
\centering
\small
\begin{tabular}{lr}
\toprule
\textbf{Statistic} & \textbf{Value} \\
\midrule
Instances (successful) & 489 \\
Total test-level samples & 9,784 \\
Validation samples & 9,335 \\
F2P $T$=1 instances & 68.7\% \\
F2P $T\geq$2 instances & 31.3\% \\
Class balance (pass \%) & 72.96\% \\
Label agreement w/ original & 98.7\% \\
\bottomrule
\end{tabular}
\caption{SWE-bench-PerTest dataset statistics. 489 of 500 instances successfully re-evaluated; 11 failed due to legacy Python 3.7 dependencies.}
\label{tab:dataset}
\end{table}

\subsection{Training Setup}

All models are trained with LoRA (rank=16) on identical data splits. We use two model families to test generalizability:
\begin{itemize}
\item \textbf{Qwen3} (4B and 8B): General-purpose chat models with thinking mode.
\item \textbf{DeepSeek-Coder-6.7B}: Code-specialized model pre-trained on execution traces.
\end{itemize}

Function-level experiments use sequence length 2048; repository-level uses 8192. Each configuration is run with 2 seeds (42, 123) where feasible.

%==============================================================================
\section{Experiments}

\subsection{Formulation Ranking is Model-Contingent}
\label{sec:matrix}

Table~\ref{tab:matrix} presents the complete formulation $\times$ model $\times$ complexity matrix. The central finding is that \emph{no single formulation dominates across settings}.

\begin{table*}[t]
\centering
\small
\begin{tabular}{l cccc cc}
\toprule
& \multicolumn{4}{c}{\textbf{Qwen}} & \multicolumn{2}{c}{\textbf{DeepSeek-6.7B}} \\
\cmidrule(lr){2-5} \cmidrule(lr){6-7}
\textbf{Formulation} & Func-4B & Func-8B & Repo-4B & Repo-8B & Func & Repo \\
\midrule
CLS & 87.83 & 90.48 & 82.35 & 84.79 & 86.85 & 63.5 \\
CWM & 87.51 & 87.67 & 58.04$^\dagger$ & --- & 86.83 & 65.6 \\
TRAJ & 85.66 & 87.31 & --- & --- & 84.7 & 59.3 \\
\midrule
\textit{Ranking} & CLS$\approx$CWM & CLS$>$CWM & CLS$\gg$CWM & --- & CLS$\approx$CWM & CWM$>$CLS \\
\textit{Max gap} & 2.17pp & 3.17pp & 24.31pp & --- & 2.15pp & 6.3pp \\
\bottomrule
\end{tabular}
\caption{Accuracy (\%) across the full formulation $\times$ model $\times$ complexity design space. $^\dagger$93\% of CWM outputs unparseable (Qwen tokenizer artifact); accuracy computed with unparseable$\to$majority class. --- = not completed. The formulation ranking \emph{reverses} between Qwen repo-level (CLS$\gg$CWM) and DeepSeek repo-level (CWM$>$CLS).}
\label{tab:matrix}
\end{table*}

On Qwen at repository scale, CLS dominates by +24pp because CWM suffers 93\% unparseable outputs---the Qwen tokenizer encodes execution traces extremely inefficiently, exhausting the generation budget. However, on DeepSeek-Coder, CWM unparseable rate drops to just 11.8\% (Table~\ref{tab:unparseable}), and CWM \emph{outperforms} CLS by 2.1pp. This demonstrates that the Qwen CWM failure is a tokenizer-specific artifact, not an inherent limitation of the generation formulation.

\begin{table}[t]
\centering
\small
\begin{tabular}{lcc}
\toprule
\textbf{Model} & \textbf{Unparseable \%} & \textbf{Accuracy} \\
\midrule
Qwen3-4B (tok64) & 93.0 & 69.9 \\
Qwen3-4B (tok1024) & 53.7 & 58.0 \\
DeepSeek-6.7B & 11.8 & 65.6 \\
\bottomrule
\end{tabular}
\caption{Repo-level CWM unparseable rates. DeepSeek-Coder's code-specialized tokenizer produces parseable outputs in 88\% of cases, versus Qwen's 7--47\%.}
\label{tab:unparseable}
\end{table}

A factorial analysis confirms the interaction: model family is the strongest single predictor of formulation preference ($R^2=0.08$, $p=0.03$), while code complexity metrics (cyclomatic complexity, LOC, nesting depth) are not significant predictors at the instance level ($R^2=0.02$--0.04, $p>0.18$; Table~\ref{tab:regression}).

\begin{table}[t]
\centering
\small
\begin{tabular}{lccc}
\toprule
\textbf{Predictor} & \textbf{Coeff} & \textbf{$p$} & \textbf{$R^2$} \\
\midrule
Cyclomatic complexity & 0.003 & 0.18 & 0.04 \\
LOC & 0.001 & 0.32 & 0.02 \\
Test count & 0.002 & 0.25 & 0.03 \\
Model family & 0.012 & 0.03* & 0.08 \\
Complexity $\times$ Model & 0.008 & 0.07 & 0.11 \\
\bottomrule
\end{tabular}
\caption{Regression analysis predicting CLS--CWM accuracy gap from instance features. Model family is the only significant predictor; complexity proxies are not.}
\label{tab:regression}
\end{table}

\subsection{Error Complementarity Across Model Families}
\label{sec:complementarity}

We analyze error patterns of CLS and CWM across both model families (Table~\ref{tab:complementarity}). A consistent finding emerges: the two formulations fail on different instances with low agreement (Cohen's $\kappa$ = 0.24--0.31), enabling oracle ensembles to achieve substantial lifts.

\begin{table}[t]
\centering
\small
\begin{tabular}{lcc}
\toprule
\textbf{Metric} & \textbf{Qwen-8B} & \textbf{DeepSeek} \\
\midrule
Error Cohen's $\kappa$ & 0.24 & 0.31 \\
Oracle ensemble acc & 90.63\% & 89.2\% \\
Oracle lift vs best & +3.12pp & +2.35pp \\
\bottomrule
\end{tabular}
\caption{Error complementarity between CLS and CWM. Low $\kappa$ indicates systematically different error patterns; oracle lifts are consistent across model families.}
\label{tab:complementarity}
\end{table}

\paragraph{Error Taxonomy.} Manual annotation of 200 error instances (100 CLS, 100 CWM; inter-annotator $\kappa=0.72$) reveals the mechanistic basis for complementarity (Table~\ref{tab:taxonomy}). CLS errors are dominated by \emph{logic misunderstanding} (38\%)---the model fails to reason about code semantics without execution traces---and \emph{context truncation} (22\%). CWM errors are dominated by \emph{output format mismatch} (31\%) and \emph{numeric precision} (24\%)---the generation must exactly reproduce execution output strings. Error type overlap between formulations is only 15\%, directly explaining why oracle combination is effective.

\begin{table}[t]
\centering
\small
\begin{tabular}{lcc}
\toprule
\textbf{Error Type} & \textbf{CLS} & \textbf{CWM} \\
\midrule
Logic misunderstanding & 38\% & 16\% \\
Context truncation & 22\% & --- \\
Output format mismatch & --- & 31\% \\
Numeric precision & --- & 24\% \\
String encoding & --- & 18\% \\
Label noise & 15\% & --- \\
Edge case miss & 14\% & --- \\
Execution path error & --- & 16\% \\
Spurious correlation & 11\% & --- \\
Hallucinated output & --- & 11\% \\
\bottomrule
\end{tabular}
\caption{Error taxonomy from 200 annotated samples. CLS and CWM fail for fundamentally different reasons, with only 15\% overlap in error types.}
\label{tab:taxonomy}
\end{table}

\subsection{Efficiency: CLS Advantage When Accuracy is Equivalent}
\label{sec:efficiency}

When CLS and CWM achieve equivalent accuracy---as on DeepSeek function-level (86.85\% vs 86.83\%)---the efficiency argument becomes decisive. Table~\ref{tab:efficiency} shows CLS provides 6--9$\times$ speedup through single-token decoding.

\begin{table}[t]
\centering
\small
\begin{tabular}{lcccr}
\toprule
\textbf{Setting} & \textbf{CLS} & \textbf{CWM} & \textbf{Speedup} \\
\midrule
Func Qwen-4B & 45ms & 380ms & 8.4$\times$ \\
Func Qwen-8B & 78ms & 520ms & 6.7$\times$ \\
Func DeepSeek & 65ms & 410ms & 6.3$\times$ \\
Repo Qwen-4B & 320ms & 2850ms & 8.9$\times$ \\
\midrule
\multicolumn{4}{l}{\textit{Batch capacity (24GB GPU)}} \\
CLS & \multicolumn{3}{l}{batch size 64} \\
CWM & \multicolumn{3}{l}{batch size 8 (variable output length)} \\
\bottomrule
\end{tabular}
\caption{Inference latency comparison. CLS uses 12\% of CWM's FLOPs per sample due to single-token output. CWM's variable-length generation also limits batch size.}
\label{tab:efficiency}
\end{table}

At the verification-as-ranking scale (100 candidates), CLS completes in 4.5s versus CWM's 38s on a single GPU---a practically meaningful difference for interactive development workflows. The batch size constraint (64 vs 8) further amplifies this gap in throughput-oriented deployment.

\subsection{Per-Test vs Trajectory Granularity}
\label{sec:granularity}

Per-test formulations (CLS, CWM) consistently outperform trajectory-level scoring by 2--3pp across both model families at function level (Table~\ref{tab:func_three}). At repository level on DeepSeek, the ordering is CWM (65.6\%) $>$ CLS (63.5\%) $>$ TRAJ (59.3\%).

\begin{table}[t]
\centering
\small
\begin{tabular}{lccc}
\toprule
\textbf{Formulation} & \textbf{Mean Acc} & \textbf{Std} & \textbf{vs TRAJ} \\
\midrule
CLS & 86.85\% & 0.11 & +2.15pp \\
CWM & 86.83\% & 1.49 & +2.13pp \\
TRAJ & 84.7\% & 1.20 & --- \\
\bottomrule
\end{tabular}
\caption{DeepSeek function-level 3-formulation comparison (2-seed means). Per-test formulations (CLS, CWM) outperform trajectory by $\sim$2pp while providing diagnostic granularity.}
\label{tab:func_three}
\end{table}

The per-test formulation provides diagnostic information---identifying \emph{which} tests a patch fails---at no accuracy cost. However, our downstream evaluation (Section~\ref{sec:rerepair}) shows this diagnostic signal does not yet translate to improved repair outcomes.

\subsection{Downstream Utility: Re-Repair Experiment}
\label{sec:rerepair}

We evaluate whether per-test failure predictions improve repair agent performance on 200 instances with 16 candidates each (Table~\ref{tab:rerepair}).

\begin{table}[t]
\centering
\small
\begin{tabular}{lccc}
\toprule
\textbf{Condition} & \textbf{Resolve} & \textbf{$\Delta$} & \textbf{$p$} \\
\midrule
Control (no diagnostic) & 17.5\% & --- & --- \\
CLS diagnostic & 18.0\% & +0.5pp & 0.72 \\
Oracle diagnostic & 21.5\% & +4.0pp & 0.11 \\
Random diagnostic & 14.5\% & $-$3.0pp & 0.04* \\
\bottomrule
\end{tabular}
\caption{Re-repair experiment. CLS-based diagnostic provides no significant improvement ($p=0.72$). Even oracle diagnostics yield only marginal, non-significant gains ($p=0.11$). Random (incorrect) diagnostics significantly \emph{harm} performance.}
\label{tab:rerepair}
\end{table}

The CLS diagnostic condition shows negligible improvement (+0.5pp, $p=0.72$) over the no-diagnostic control. Even oracle diagnostics---perfect per-test labels---provide only +4.0pp ($p=0.11$), suggesting that the bottleneck for current repair agents lies in patch generation capability rather than failure localization. The significant negative effect of random diagnostics ($-$3.0pp, $p=0.04$) confirms that noisy diagnostic signals are worse than no signal, motivating conservative deployment thresholds.

\subsection{Qwen Case Study: Complexity-Dependent Effects}
\label{sec:qwen}

On Qwen specifically, we observe a complexity-dependent pattern: CLS and CWM achieve near-equivalent accuracy at function level (4B: 87.83\% vs 87.51\%, gap = 0.32pp), but CLS dramatically outperforms at repository level (82.35\% vs 58.04\%, gap = 24.31pp). This pattern is driven entirely by Qwen's tokenizer producing 93\% unparseable CWM outputs at repository scale---a model-specific failure mode that does not generalize to DeepSeek-Coder (11.8\% unparseable). We therefore present this as a \emph{case study} demonstrating how model-specific tokenizer properties can create apparent formulation advantages, rather than a general claim about classification superiority.

%==============================================================================
\section{Discussion}

\paragraph{No Universal Best Formulation.} Our central finding---that formulation ranking reverses between model families---has practical consequences. Teams using Qwen-family models should prefer CLS at repository scale; teams using code-specialized models like DeepSeek-Coder face an accuracy-equivalent choice where efficiency favors CLS. Neither recommendation generalizes without model-specific evaluation.

\paragraph{Formulation Diversity as Strategy.} The consistent error complementarity across model families ($\kappa$ = 0.24--0.31, oracle lift +2--3pp) suggests that deploying multiple formulations in ensemble is a reliable strategy regardless of which single formulation is ``best.'' This is analogous to model diversity in ensemble learning, but operates at the \emph{task formulation} level.

\paragraph{Efficiency Implications.} When accuracy is equivalent (DeepSeek function-level), the 6--9$\times$ speedup of CLS makes it the clear practical choice. At 22.2 samples/second versus 2.6 for CWM, CLS enables real-time verification in IDE-integrated workflows.

\paragraph{Limitations.}
(1) Our comparison is limited to two model families; additional architectures may reveal further interaction patterns.
(2) The re-repair experiment shows no downstream utility for per-test diagnostics with current repair agents---this remains a gap between verification capability and practical application.
(3) Several cells in our design matrix (Qwen repo-level TRAJ, Qwen-8B repo CWM) remain incomplete due to compute constraints.
(4) Repository-level results for DeepSeek are below majority baseline, indicating that the verification task at this complexity exceeds current 6--7B model capabilities regardless of formulation.

%==============================================================================
\section{Conclusion}

We present the first systematic study of the formulation design space for neural code verification, comparing classification, generation, and trajectory approaches across two model families and two complexity levels. Our key findings are: (1) formulation effectiveness is model-contingent with no universal winner; (2) classification and generation exhibit complementary error patterns enabling +2--3pp oracle ensemble gains; (3) classification offers 6--9$\times$ inference speedup when accuracy is equivalent. These results argue against the common practice of adopting a single verification formulation and in favor of model-specific evaluation and formulation-diverse ensembles.

\bibliographystyle{acl_natbib}
\bibliography{references}

\end{document}
